\begin{lemma}[Case 2, observation (19)]
Let \(G\) be a finite loopless multigraph, let \(\varrho_1\) be a
partial \(k\)-edge-coloring of \(G\), and let \(T\) and \(T'\) be
arbitrary tree-sequences for \(G\) with the same distinguished first
edge.  Let \(S\) be an arbitrary finite set of colors, let \(v_*\) be a
distinguished vertex of \(T'\), and let \(\delta\) be a color.  Suppose
that \(T'\) is a \(\varrho_1\)-closure of \(T\), that \(v_*\in V(T')\),
\[
 \overline{\varrho}_1(v)\cap S=\varnothing
 \quad\text{for every }v\in V(T')\setminus\{v_*\},          \tag{19a}
\]
and
\[
              \overline{\varrho}_1(v_*)\cap S=\{\delta\}.    \tag{19b}
\]
Then
\[
             \overline{\varrho}_1(T')\cap S=\{\delta\},      \tag{19c}
\]
and every exact tree edge appended after \(T\) in the closure sequence is
not assigned any color in \(S\setminus\{\delta\}\).  More explicitly,
for every index \(a\) with
\(0\leq a<\operatorname{length}(T')\) and
\(\operatorname{length}(T)\leq a\), and every
\(c\in S\setminus\{\delta\}\),
\[
               \varrho_1\bigl(T'.\operatorname{edge}(a)\bigr)
                    \ne \operatorname{some}(c).              \tag{19d}
\]
\end{lemma}
\begin{proof}
By \(\noderef{TreeMissingColors}\), membership of a color \(c\) in
\(\overline{\varrho}_1(T')\) is witnessed by a displayed vertex
\(v\in V(T')\) at which \(c\) is missing.  Suppose also that \(c\in S\).
If \(v=v_*\), equation (19b) gives \(c=\delta\).  If \(v\ne v_*\),
then the definition of \(\noderef{TreeSequenceVertexSet}\) puts \(v\) in
\(V(T')\setminus\{v_*\}\), and (19a) contradicts the simultaneous
membership of \(c\) in the missing set of \(v\) and in \(S\).  This proves
the inclusion from left to right in (19c).  Conversely, (19b) says that
\(\delta\) is missing at \(v_*\), and \(v_*\in V(T')\); hence
\(\delta\in\overline{\varrho}_1(T')\) by
\(\noderef{TreeMissingColors}\), while (19b) also puts it in \(S\).
Thus (19c) is an equality.

It remains to retain the identities and order of the closure edges.
Unfold \(\noderef{TashkinovClosure}\).  Its first clause, together with
\(\noderef{TashkinovAugmentationSequence}\), supplies a finite chain
\[
                         T^0=T,T^1,\ldots,T^m=T'.             \tag{1}
\]
of \(\varrho_1\)-Tashkinov augmentations.  Unfolding
\(\noderef{TashkinovAugmentation}\) at step \(r<m\) gives all of the
following exact information: \(T^{r+1}\) has length
\(|T^r|+1\), \(T^r\) is an ordered initial segment of \(T^{r+1}\), and
the new edge identity at position \(|T^r|\) has a color \(d_r\) which is
missing on \(T^r\).

Put \(q=\operatorname{length}(T)\).  Induction on \(r\) gives
\(|T^r|=q+r\).  Composition of the length and
entrywise vertex and edge equalities in
\(\noderef{TreeSequenceInitialSegment}\) shows that each \(T^r\) is an
ordered initial segment of \(T'\).  In particular,
\[
                         V(T^r)\subseteq V(T'),               \tag{2}
\]
and the exact edge appended at step \(r\) is still the exact displayed
edge of \(T'\) at position \(q+r\); parallel edges are not identified
in this comparison.

Now fix a displayed edge position \(a\ge q\) of \(T'\).  Since
\(a<|T'|=q+m\), the natural number \(r=a-q\) satisfies \(r<m\), and
the edge at position \(a\) is precisely the edge appended in the passage
from \(T^r\) to \(T^{r+1}\).  Let its \(\varrho_1\)-color be \(d_r\).
The augmentation certificate says
\(d_r\in\overline{\varrho}_1(T^r)\).  By (2), the vertex witness for
this missing color also belongs to \(V(T')\); the union definition in
\(\noderef{TreeMissingColors}\) therefore yields
\[
                         d_r\in\overline{\varrho}_1(T').      \tag{3}
\]
If \(d_r\in S\setminus\{\delta\}\), then (19c) and (3) force
\(d_r=\delta\), while membership in the displayed set difference says
\(d_r\ne\delta\), a contradiction.  Thus the exact appended edge at
position \(a\) is not assigned any color in \(S\setminus\{\delta\}\).
Since \(a\) was arbitrary, this is (19d).

The manuscript's observation (19) is the specialization of this generic
lemma obtained when the later Case 2 data take
\(T=T_j(v_n)\), \(T'=T'_j\), \(v_*=v_n\), and
\(S=S'=\bigcup_{i\in L_1}S_i\).  Its Case 2 hypotheses are exactly
(19a) and (19b), so the two conclusions above give the displayed
observation there.
\end{proof}
